\subsection{Learning a representation versus exploiting its structure}

OPQ learns rotation and subspace codebooks to reduce distortion
\cite{opq}; JHQ's random square rotation supports an analytical primary
model. A learned replacement can improve reconstruction while losing the
Cartesian structure needed for exact 16+16 splitting.
ScaNN weights quantisation error for maximum-inner-product retrieval
\cite{scann}, illustrating why low reconstruction MSE need not preserve
nearest-neighbour ranking. Chapter 8 tests that distinction within JHQ.

\subsection{Small lookup tables have substantial prior art}

PQ Fast Scan targets CPU SIMD-register tables \cite{fastscan};
Quicker ADC extends register lookup with irregular quantisers and split
tables \cite{quickeradc}. These are precedents for co-designing table
shape and execution, not GPU throughput measurements. JHQ's exact
Cartesian split still requires numerical and cutoff-tie checks.

RaBitQ combines randomised compact codes with an analysed distance
estimator \cite{rabitq}; multi-bit extensions trade storage for accuracy
\cite{rabitq_multi}. JHQ instead uses signs for primary retention and
stored residuals for refinement, so RaBitQ's guarantees do not transfer.
GPU IVF-RaBitQ supplies GPU construction and fused search
\cite{ivfrabitq}, making it a contemporary end-to-end baseline. E1 tests
multiple bit budgets and reveals high-recall reversals.

\subsection{Refinement effort and adaptive stopping}

Faiss's \texttt{k\_factor} and cuVS expose candidate expansion/refinement
\cite{faiss_refine,cuvs_refinement}. Faiss autotuning uses GT-labelled
queries \cite{faiss_autotune}; learned termination predicts effort
\cite{li_termination}, while DARTH targets declarative recall
\cite{darth}. Chapter 6 calibrates one unlabelled workload-level multiplier
for reuse. It avoids a learned predictor but inherits reference error.
